\section{Introduction}
\label{sec:intro}

Wildlife management runs on counts. Deer density drives hunting quotas, culling decisions,
disease-response planning, and vehicle-collision mitigation, and those counts are still
produced largely by human observers riding transects at night with spotlights or handheld
thermal scopes. Automating the count is an obvious application of modern detection, and the
literature has largely treated it as one: train a detector on annotated animal imagery,
report mAP, and declare the survey problem addressed.

This paper argues, and then measures, that the two problems come apart. A detector is
rewarded for finding animals in \emph{frames}. A survey needs the number of distinct
\emph{animals}. Between those sits a chain of decisions (associating detections into tracks, deciding which tracks are real, deciding which are duplicates of an animal already counted) that no detection metric scores and that, we find, dominates the error.

That distinction also fixes the yardstick we use throughout. A survey does not need a box
placed precisely on an animal; it needs the animal noticed once. We therefore score the
detection stage \emph{per animal}, over the whole of an animal's track, asking whether any
frame produced a box touching it, and report it alongside the conventional per-box metrics
rather than in place of them. The two measure different things, and that difference is the
paper's starting point: every one of the twelve architectures we train finds between 82\% and
99.6\% of the 235 annotated individuals in at least one frame, while none exceeds
AP$_{50}$ 0.510 on the same images. Per-box mAP remains a valid measure of localization. It
answers a question the survey does not ask, and read as a survey metric it would place the
difficulty in detection, where these measurements show it is very nearly solved. That claim is
bounded, and we bound it ourselves: it holds within this corpus, where the detector has seen
every site. Under leave-one-site-out training it weakens to $90.7\%$ of animals reached
(\cref{sec:loso}), which is still high but no longer near-saturated.

\begin{figure}[H]
  \centering
  \includegraphics[width=\linewidth]{figures/fig1_funnel.pdf}
  \caption{\textbf{Detection is not the bottleneck.} Of the 83 deer in the 13 videos the
  detector never trained on, up to 82 are found by \emph{some} candidate track and 73
  acquire a distinct track of their own, but only 52--55 are counted. Every bar is one of
  four candidate-generation configurations (\cref{sec:inversion}); improving the pool moves
  the blue and green bars up and the orange bar \emph{down}.\label{fig:funnel}}
\end{figure}

We study this on a corpus we built for the purpose: 32 FLIR thermal road transects from
four sites in Illinois, 521,930 frames at $640\times512$ and 60\,fps, with every individual
deer annotated as a distinct track, 236 animals, 21,646 boxes. Because one annotated
track is one animal, the corpus supplies detection, tracking and count ground truth
simultaneously, the property on which the decomposition below depends.

Our central instrument is a three-stage decomposition of counting error. For each
ground-truth animal we ask whether it is \reachedq (touched by at least one candidate
track), whether it is \primaryq (some candidate covers it better than it covers any other
animal, so a one-count-per-track decision \emph{could} count it), and whether it is finally
\countedq. The gaps between these three numbers attribute error to detection, to
association, and to confirmation respectively. \Cref{fig:funnel} gives the result: on
held-out video the pipeline reaches $98.8 \pm 5.3\%$ of animals, gives $88.0 \pm 8.7\%$ their
own track, and counts $62.7 \pm 10.7\%$ (95\% Wilson intervals over the 83 animals).

Having localized the loss, we attack it. We improve candidate generation four independent
ways, lowering detector confidence from 0.10 to 0.02, removing the tracker's
track-initialization gate, raising the NMS IoU threshold from 0.50 to 0.90, and adding
appearance re-identification to the association step. Each one demonstrably improves the
pool. Each one makes the final count worse or no better. We then replace the hand-tuned
confirmation rule with learned alternatives of varying capacity, under three evaluation
protocols, and every one of them loses to three hand-tuned parameters.

Neither outcome is a gap in the search for a method; together they are the finding. The
confirmation stage is bounded by supervision rather than by capacity: roughly 200 positive
tracks against 7,000--28,000 negatives, with the positives not linearly separable from
the duplicates in any feature we measured. Growing the candidate pool grows the negatives
far faster than the positives, which is exactly why a richer pool returns a lower count.

This paper contributes three things. The first is the corpus: 32 fully annotated thermal
road transects holding 236 individually tracked deer, with per-individual identity that makes
detection, tracking and counting evaluable against one set of annotations, released with the
evaluation harness. The second is the decomposition above, together with a held-out protocol
that removes the optimism of in-sample evaluation: reporting counting accuracy on videos the detector trained on inflates it by 15.5 points. The third is the pair of negative results
the decomposition exposes: counting accuracy moves opposite to candidate quality across four
independent interventions, and learned confirmation loses to a three-parameter rule under
three evaluation protocols. The rest of the paper (the twelve-detector benchmark, the matching-criterion analysis, the sensitivity of the confirmation rule, the failure of appearance re-identification at 27 pixels, and what survives a change of site) is evidence
for those three or a consequence of them.
